%AttackName: Bullseye Polytope

%Input:A set of base images from the intended misclassification class (e.g., 'ship'), a target image (or images) from the target class (e.g., 'frog'), a pre-trained feature extractor (or substitute networks), and a perturbation budget ε.
%Output:A set of clean-label poison samples, each visually similar to its base image but imperceptibly perturbed, such that their feature representations form a convex polytope whose center is close to the target's feature vector. When these poisons are injected into the training set, a model fine-tuned on this data will misclassify the target as the poison class.
%Formula:minimize_{x_p^{(j)}} (1/(2m)) * sum_{i=1}^m || φ^{(i)}(x_t) - (1/k) * sum_{j=1}^k φ^{(i)}(x_p^{(j)}) ||^2 / ||φ^{(i)}(x_t)||^2, subject to ||x_p^{(j)} - x_b^{(j)}||_∞ ≤ ε, ∀j
%Explanation:Bullseye Polytope is a clean-label data poisoning attack targeting transfer learning. The attacker crafts poison samples by imperceptibly perturbing base images from the poison class so that, in the feature space of one or more substitute networks, the mean feature vector of the poisons is as close as possible to the target's feature vector. This is achieved by fixing the convex coefficients to be uniform (i.e., the target is at the center of the poisons' convex hull), which improves both the transferability and speed of the attack compared to prior work. When these poisons are injected into the victim's training set, a model fine-tuned on this data is likely to misclassify the target as the poison class, even in black-box or gray-box settings.